\exercises

\tierunderstand

\begin{bt}
  \item Chặn~\eqref{eq:ch03-chan-tung-thua-so} dùng bất đẳng thức chuẩn của tích
    không vượt quá tích các chuẩn. Viết lại bước quy nạp dẫn
    tới~\eqref{eq:ch03-chan-mu}, và chỉ ra chính xác chỗ nào cần
    \(\eta < 1\) chứ không phải \(\eta \le 1\).

  \item Với sigmoid, \(\gamma = 1/4\), nên ngưỡng của \tn{điều kiện đủ}{sufficient condition} là 4 chứ
    không phải 1. Giải thích vì sao điều đó không có nghĩa là mạng dùng sigmoid
    nhớ được lâu hơn mạng dùng tanh.

  \item Cho \(\mat{W} = \begin{psmallmatrix} 0.9 & 3 \\ 0 & 0.9 \end{psmallmatrix}\).
    Viết \(\mat{W}^{\,l}\) ở dạng đóng, rồi tìm \(l\) làm phần tử ngoài đường
    chéo đạt cực đại. So kết quả với đỉnh đo được ở
    mục~\ref{sec:ch03-khe-ho}.

  \item Mục~\ref{sec:ch03-khe-ho} nói phát biểu của bài đúng ở nghĩa tiệm cận và
    sai ở nghĩa từng bước. Phát biểu lại điều kiện đủ cho vanishing sao cho nó
    đúng ở cả hai nghĩa, và nói rõ giả thiết bạn phải thêm vào.

  \item Nhìn hình~\ref{fig:ch03-vach-mat-loi}. Nếu learning rate được chọn đủ
    nhỏ để một bước tại đỉnh vách không vượt ra ngoài, bước đó dài bao nhiêu ở
    vùng phẳng? Dùng hai số đo trong mục~\ref{sec:ch03-he-dong-luc} để trả lời,
    rồi nói xem con số ấy có dùng được để huấn luyện không.

  \item Mục~\ref{sec:ch03-con-lai-gi} lập luận rằng hai biện pháp của bài phải
    đi cùng nhau. Dựng lại lập luận đó bằng ngôn ngữ \tn{miền hút}{basin of attraction} và ranh giới miền
    hút, không dùng lại câu chữ của sách.
\end{bt}

\tierapply

Trạng thái xuất phát cho cả bốn bài: clone repo companion, đứng ở thư mục gốc
của nó, và checkout đúng tag của chương này.

\begin{minted}{console}
$ git clone https://github.com/Giang-Dang/rnn-to-transformer-lab.git
$ cd rnn-to-transformer-lab
$ git checkout ch03
$ conda env create -f environment.yml
$ conda activate rnn-to-transformer-lab
$ python verify.py --only ch03
\end{minted}

Dòng cuối phải in ra \code{verify: ok}. Nếu không, đừng làm tiếp: mọi con số
bạn đo sau đó sẽ không so được với sách.

\begin{bt}
  \item \code{experiments/ch03\_decay.py} chạy với tanh. Đổi sang sigmoid và
    chạy lại bốn \tn{bán kính phổ}{spectral radius} cũ. Dự đoán trước khi chạy: dòng nào đổi nhiều
    nhất, và theo chiều nào? Đối chiếu với \(\gamma = 1/4\).

  \item \code{jordan\_block(2, 0.9, 3.0)} tách bán kính phổ khỏi \tn{chuẩn phổ}{spectral norm} bằng
    một phần tử ngoài đường chéo. Viết một hàm dựng ma trận \(n \times n\) ngẫu
    nhiên có bán kính phổ đặt trước nhưng chuẩn phổ lớn hơn một hệ số cho trước,
    rồi đo lại đỉnh và bước đầu tiên tụt xuống dưới 1. Kiểm rằng hàm của bạn
    không vô tình dựng ra \tn{ma trận chuẩn tắc}{normal matrix}.

  \item \code{experiments/ch03\_clipping.py} đi đúng một bước. Sửa nó thành một
    vòng lặp 50 bước từ cùng điểm xuất phát, chạy ba lần với ngưỡng
    \code{None}, \code{1.0} và \code{0.1}, và in chi phí sau mỗi bước. Ngưỡng
    nào xuống thấp nhất sau 50 bước, và ngưỡng ấy có phải ngưỡng bạn đoán
    không?

  \item Ngưỡng trong \code{clipping.py} so bằng \code{>=}. Đổi thành \code{>} và
    chạy \code{python -m pytest -q}. Test nào hỏng, và nếu không có test nào
    hỏng thì điều đó nói gì về bộ test? Viết thêm một test bắt được khác biệt
    ấy, hoặc lập luận rằng khác biệt ấy không quan sát được và giải thích vì
    sao.
\end{bt}

\tierextend

\begin{bt}
  \item Mọi ma trận hồi quy trong chương này đều do tôi đặt phổ bằng tay. Lấy
    một mạng thật đã huấn luyện xong, đo \(\rho(\mat{W}_{hh})\) và
    \(\norm{\mat{W}_{hh}}\) sau mỗi epoch, và xem khoảng cách giữa hai số ấy
    thay đổi ra sao trong lúc huấn luyện. Huấn luyện đẩy ma trận về phía chuẩn
    tắc hay ra xa nó?

  \item Bài đưa heuristic chọn ngưỡng: nhìn chuẩn gradient trung bình qua nhiều
    bước. Mười ba năm sau, ngưỡng clipping vẫn là một siêu tham số phải dò cho
    từng tác vụ, và chưa có quy tắc nào độc lập tác vụ được chấp nhận rộng rãi.
    Đây là một chỗ còn mở thật sự. Thử đề xuất một quy tắc dựa trên đại lượng
    bạn đo được trong lúc huấn luyện, rồi tìm một tác vụ làm quy tắc đó hỏng.

  \item Tăng trưởng tạm thời ở mục~\ref{sec:ch03-khe-ho} là hệ quả của tính
    không chuẩn tắc. Kiến trúc nào trong các chương sau còn giữ một ma trận
    được nhân lặp lại nhiều lần như vậy, và kiến trúc nào thì không? Trả lời
    được câu này trước khi đọc chương~\ref{ch:transformer} sẽ làm chương đó dễ
    hơn nhiều.

  \item \(\Omega\) trong~\eqref{eq:ch03-omega} phạt mọi bước lệch khỏi tỷ số 1,
    kể cả bước lớn hơn 1. Thiết kế một biến thể chỉ phạt phía co lại, đo nó trên
    cùng ba mạng ở mục~\ref{sec:ch03-regularizer}, và giải thích vì sao bài báo
    có thể đã chọn dạng đối xứng dù dạng một phía nghe hợp lý hơn.
\end{bt}
